% §17.1 What is a Cipher Map?
%
% Cross-reference resolution status:
%
% DEFINED here:
%   \label{sec:what-is-cipher-map}            -- this section
%   \label{def:cipher-map}                    -- the cipher-map definition
%
% ALREADY DEFINED (chapter 1):
%   \Cref{sec:noisy-bit}                      -- chapter 1 §1.2 fprate/fnrate
%
% ALREADY DEFINED (chapter 5):
%   \Cref{sec:bloom-from-scratch}             -- chapter 5 §5.2 Bloom filter
%
% FORWARD refs within chapter 17:
%   \Cref{sec:bernoulli-error-model}          -- §17.2 (Task 3)
%   \Cref{sec:fhe-sidebar}                    -- §17.4 (Task 5)
%
% §17.1 introduces the cipher-map abstraction: a function that computes
% faithfully on encrypted inputs up to bounded false-positive and
% false-negative rates. Five items: (1) plaintext setup; (2) cipher-map
% definition; (3) why approximate is acceptable; (4) the trapdoor key;
% (5) three example cipher maps.

\section{What is a Cipher Map?}
\label{sec:what-is-cipher-map}

The first sixteen chapters built a framework for noisy computation on
unencrypted data. The Bloom filter of chapter 5 returns wrong answers
because the structure is small; the classifier of chapter 9 returns
wrong answers because the world is hard; the sketch of chapter 10
returns wrong answers because the stream is large. In every case the
error is a property of the structure, and the user accepts it for
reasons of space or time.

This chapter takes the same Bernoulli machinery and applies it to a
different setting. The data is encrypted; the computation runs on the
encrypted form; only a holder of a secret key can read the result.
Bloom filters and noisy classifiers reappear, but now their false
positives and false negatives serve a second purpose alongside their
operational role: the bounded error is also what keeps the underlying
plaintext hidden from the party doing the computation.

\subsection*{The plaintext setting}

Begin with a function on plaintext data. Let $X$ be a domain of
plaintext values (records, sets, scalars, whatever the application
requires), let $Y$ be a codomain, and let
\[
  f : X \to Y
\]
be the function the user wants to compute. In the cleartext setting
the user holds an $x \in X$, applies $f$, and reads $f(x)$. Chapters
1 through 16 sometimes replaced $f$ with an approximate $\tilde f$
that occasionally returned the wrong $y$; the error analysis was
operational.

The cipher-map setting introduces a different actor. A holder of $x$
(the \emph{client}) wants $f(x)$ but does not want to give $x$ to the
machine doing the computation. Instead, the client encrypts $x$ under
an encryption map $\mathsf{Enc}$, sends the ciphertext to a
\emph{server}, and asks the server to compute something on
$\mathsf{Enc}(x)$ that the client can decrypt to obtain (something
close to) $f(x)$. The server never sees $x$. The server also never
sees $f(x)$ directly; what the server returns is itself an encryption
of an output, which the client decrypts using $\mathsf{Dec}$.

The shape of the computation is therefore not $f$ itself but a
companion map that operates on the encrypted representations. That
companion map is the cipher map.

\subsection*{The cipher-map abstraction}

\begin{definition}[Cipher map]
\label{def:cipher-map}
Let $f : X \to Y$ be a plaintext function. Let $\mathsf{Enc} : X \to
\mathsf{Enc}(X)$ be an encryption map onto an encrypted-domain space
$\mathsf{Enc}(X)$, with corresponding decryption $\mathsf{Dec} :
\mathsf{Enc}(Y) \to Y$. A \emph{cipher map} for $f$ is a function
\[
  \tilde f : \mathsf{Enc}(X) \to \mathsf{Enc}(Y)
\]
on the encrypted representations such that, for $x \in X$,
\[
  \Pr\left[\,\mathsf{Dec}(\tilde f(\mathsf{Enc}(x))) = f(x)\,\right]
  \;\ge\; 1 - \fnrate
\]
when $f(x)$ is a real positive output, and
\[
  \Pr\left[\,\mathsf{Dec}(\tilde f(\mathsf{Enc}(x))) = y \text{ for some
  spurious } y \ne f(x)\,\right] \;\le\; \fprate
\]
when $f(x)$ is the negative case. The pair $(\fprate, \fnrate)$ is
the cipher map's \emph{error profile}, in the Bernoulli sense of
chapter 1 \S\ref{sec:noisy-bit}.
\end{definition}

Three pieces of the definition deserve emphasis. The cipher map
operates entirely on the encrypted domain: every step from
$\mathsf{Enc}(x)$ through $\tilde f$ to $\mathsf{Enc}(y)$ leaves the
ciphertext encrypted. The decryption happens on the client side,
after the encrypted output has been returned. And the error rates
$\fprate, \fnrate$ are exactly the rates that the framework has
spent sixteen chapters analyzing: there is no new probabilistic
machinery to introduce; the cipher-map error is a Bernoulli error.

What is new is the second meaning carried by the error. A cipher map
with $\fprate = 0$ and $\fnrate = 0$ behaves identically to $f$ from
the client's view, but it also reveals more information about $x$ to
any adversary that observes $\tilde f$'s behavior across many
queries. A cipher map with positive error obscures $x$ in proportion
to the error magnitude. \Cref{sec:bernoulli-error-model} makes this
duality precise; \Cref{sec:trapdoor-paradigm} unpacks its
operational consequences.

\subsection*{Why approximate is acceptable}

It is reasonable to ask why one would accept bounded error in a
cryptographic setting. The cleartext approximate-set literature
tolerates Bloom-filter false positives because the alternative
(storing the full set with exact membership) costs too much memory;
the cipher-map literature tolerates cipher-map false positives for
an analogous reason in the cryptographic dimension.

The exact alternative exists. Fully homomorphic encryption (FHE),
introduced by Gentry in 2009 and refined by the BFV, BGV, and CKKS
schemes since, provides an $\tilde f$ with $\fprate = \fnrate = 0$
for any computable $f$. FHE is the cryptographer's answer to the
question ``can we compute exactly on encrypted data?'' and the
answer is yes. The catch is that FHE ciphertexts are large
(kilobytes to megabytes per encrypted bit in practical schemes) and
slow (seconds per encrypted bit operation). For workloads where the
plaintext computation runs in microseconds, an FHE realization can
run six orders of magnitude slower with three orders of magnitude
more memory.

Cipher maps trade error for practicality. A Bloom-filter-based
cipher map for set membership uses on the order of
$-\log_2 \fprate + \mu$ bits per inserted element (the optimal
formula of chapter 14), runs in constant time per query, and gives
a false-positive rate the client can tune. The trade-off is the
familiar Bernoulli one: space and runtime are exchanged for
controlled error. \Cref{sec:fhe-sidebar} returns to the
FHE-versus-cipher-map comparison in detail.

\subsection*{The trapdoor key}

The key that distinguishes a cipher map from any other noisy
function is the \emph{trapdoor key}: a secret that enables recovery
of $f(x)$ from $\tilde f(\mathsf{Enc}(x))$. The client holds the
trapdoor; the server does not. Without the trapdoor, the server
sees only opaque bit patterns flowing through $\tilde f$ and cannot
distinguish meaningful queries from filler, real outputs from noise,
or one plaintext value from another. With the trapdoor, the client
can decrypt the cipher-map output and recover the answer to within
the cipher map's $\fprate, \fnrate$ bounds.

Three properties make a trapdoor key worth its name. First, the
trapdoor is short: a key of constant length, not a representation
that grows with the data. Second, the trapdoor is necessary: no
polynomial-time computation without it can recover plaintext from
the cipher-map output. Third, the trapdoor is sufficient: a holder
of the key can decrypt every cipher-map output with the same per-call
cost as a plaintext function evaluation. These three properties
distinguish a cipher map from a generic encryption (which lacks
property three) and from a generic noisy function (which lacks
properties one and two).

\subsection*{Three example cipher maps}

Three cipher maps surface repeatedly in this chapter and in
chapter 18. Each is sketched here at the level needed to ground
\Cref{sec:bernoulli-error-model}'s analysis.

\textbf{Encrypted set membership.} Let $A \subset X$ be a plaintext
set the client wants to query without revealing to the server. The
function is the indicator $f_A(x) = \mathbf{1}[x \in A]$. The
client builds a Bloom filter on $A$ using a keyed hash family (the
key is the trapdoor), encrypts the resulting bit array under a
suitable scheme, and gives it to the server. To query, the client
encrypts the query element $x$, sends $\mathsf{Enc}(x)$ to the
server, and asks the server to test the encrypted bits. The server
returns a (possibly encrypted) bit. With probability at least
$1 - \fprate$ the bit is correct on negative queries
($x \notin A$); the bit is always correct on positive queries
($\fnrate = 0$), reflecting the one-sided error of the Bloom
construction.

\textbf{Encrypted equality test.} Two parties want to know whether
$x = y$ without revealing $x$ or $y$. The function is
$f_=(x, y) = \mathbf{1}[x = y]$. Under deterministic encryption,
identical plaintexts encrypt to identical ciphertexts, so equality
on ciphertexts coincides with equality on plaintexts and the cipher
map has $\fprate = \fnrate = 0$. The cost is a known information
leak: an adversary observing many equality tests can build the
distribution of plaintext values from ciphertext frequencies. This
is a cipher map with no Bernoulli error but with non-zero
information leak: a useful comparison case, since the framework's
$\fprate, \fnrate$ are not the only confidentiality knob.

\textbf{Encrypted ordering.} A client wants to know whether
$x \le y$ on encrypted inputs (for sorting, ranged queries, or
threshold tests). The function is $f_{\le}(x, y) =
\mathbf{1}[x \le y]$. Under order-preserving encryption, the
ciphertext encoding preserves the plaintext order, so the server
can compare ciphertexts directly. The cipher map has $\fprate =
\fnrate = 0$ on the comparison itself but leaks the entire
ranking of the underlying plaintexts to the server: an even larger
leak than the equality case. Order-preserving encryption is again
a cipher map with no Bernoulli error and a tightly characterized
leakage profile.

These three examples span the rate landscape. The Bloom-based
cipher map sits at the high-$\fprate$, zero-$\fnrate$ corner of the
perfect-filter line of chapter 6; the equality and ordering cipher
maps sit at the zero-zero corner but pay their confidentiality cost
through structural leakage rather than rate. The chapter's main
work is to show that the framework's $\fprate, \fnrate$ machinery
gives a quantitative grip on the first case, and that the framework's
PPV/NPV machinery (chapter 7 §7.4) gives a confidentiality measure
that includes the second.
